\chapter{CONCEPTION DE HAUT NIVEAU}
\thispagestyle{empty}

La conception de haut niveau du système se compose de quatre couches principales qui s'alimentent mutuellement : Préparation des Données, Encodage Textuel, Apprentissage sur Graphe et Présentation des Recommandations.

% TODO: ajouter un diagramme ici
% \begin{figure}[H]
%     \centering
%     \includegraphics[width=\textwidth]{architecture.png}
%     \caption{Architecture de haut niveau du système.}
%     \label{fig:architecture}
% \end{figure}

\section{Couche 1 -- Préparation des données}

Cette couche est responsable de la transformation des données brutes de Wikipédia en une structure adaptée à l'entraînement du modèle.

\begin{enumerate}
    \item La topologie de graphe (matrice d'adjacence) du jeu de données Wiki-CS est extraite.
    \item Les titres des n\oe uds sont normalisés et mis en correspondance avec les titres réels des articles sur Wikipédia.
    \item Les textes des articles sont obtenus en priorité à partir du Wikipedia STEM Corpus, et ceux non trouvés sont complétés via l'API Wikipédia.
    \item Les textes obtenus sont nettoyés et transmis à la couche suivante.
\end{enumerate}

Sortie : un graphe orienté de 11 701 n\oe uds (matrice d'adjacence) et le texte de l'article correspondant à chaque n\oe ud.

\section{Couche 2 -- Encodage textuel}

Dans cette couche, le texte de chaque article est transformé en vecteurs denses via des modèles pré-entraînés de plongements de phrases. Deux modèles sont évalués de manière comparative :

\begin{itemize}
    \item BAAI/bge-m3 $\rightarrow$ vecteur de plongement de 1024 dimensions
    \item Alibaba-NLP/gte-modernbert-base $\rightarrow$ vecteur de plongement de 768 dimensions
\end{itemize}

Ces vecteurs seront utilisés comme attributs de n\oe uds (node features) sur le graphe.

Sortie : matrice d'attributs des n\oe uds $\mathbf{X} \in \mathbb{R}^{N \times d}$ (où $N = 11\,701$, $d \in \{768, 1024\}$).

\section{Couche 3 -- Apprentissage sur graphe (GAT)}

Cette couche combine les vecteurs de plongement textuel avec la structure de graphe pour apprendre des représentations enrichies des n\oe uds.

\begin{enumerate}
    \item La matrice d'attributs des n\oe uds $\mathbf{X}$ et la matrice d'adjacence $\mathbf{A}$ sont fournies en entrée au modèle GAT.
    \item Le mécanisme d'attention multi-têtes du GAT calcule le poids d'importance pour chaque arête.
    \item Plusieurs couches GAT sont appliquées séquentiellement, permettant à chaque n\oe ud d'obtenir une représentation enrichie intégrant l'information de voisinage multi-sauts.
    \item Le GCN est entraîné séparément sur les mêmes données comme modèle baseline, et les résultats sont comparés.
\end{enumerate}

Sortie : matrice de plongements enrichis des n\oe uds $\mathbf{Z} \in \mathbb{R}^{N \times d'}$, où chaque vecteur reflète à la fois le contenu sémantique de l'article et sa position dans la structure du graphe.

\section{Couche 4 -- Recommandation et présentation hiérarchique}

Cette couche prend la requête de l'utilisateur et produit la recommandation finale.

\begin{enumerate}
    \item La requête de l'utilisateur est vectorisée avec le même modèle de plongements de phrases.
    \item La similarité entre le vecteur de requête et les vecteurs de plongement des n\oe uds ($\mathbf{Z}$) est calculée (similarité cosinus) et les $k$ articles les plus proches sont identifiés comme candidats.
    \item Les articles candidats sont transmis à un LLM exécuté localement (DeepSeek-R1 ou similaire).
    \item Le LLM organise les articles candidats en un parcours d'apprentissage hiérarchique selon les relations conceptuelles de prérequis et les présente à l'utilisateur.
\end{enumerate}

Sortie : un parcours structuré d'apprentissage d'articles présenté à l'utilisateur, répondant à la question ``par où commencer ?''.